\section{ماژول flap (Fluctuation-Based Adaptive Structured Pruning)}
\label{sec:module_flap}

\subsection{معرفی و انگیزه}
هرس ساختاریافته سنتی در مدل‌های ترنسفورمر اغلب با اتکا بر وزن‌های ایستا (مانند نُرم $\ell_1$ یا $\ell_2$) انجام می‌شود که بدون در نظر گرفتن دینامیک فعال‌سازی‌های ورودی، ستون‌ها یا کانال‌هایی را حذف می‌کند که در واقعیت برای حفظ بازنمایی داده‌ها حیاتی هستند. از سوی دیگر، روش‌های متکی بر اطلاعات مرتبه دوم، پیچیدگی محاسباتی بالایی داشته و سربار زمانی قابل توجهی در پردازش شبکه‌های بزرگ ایجاد می‌کنند.

ماژول فشرده‌سازی \lr{FLAP (Fluctuation-Based Adaptive Structured Pruning)} به عنوان یک راهکار هرس ساختاریافته تطبیقی مبتنی بر نوسانات فعال‌سازی پدید آمده است. ایده محوری این روش، بهره‌گیری از معیار واریانس فعال‌سازی‌های ورودی (\lr{Activation Fluctuation/Variance}) به عنوان عاملی کلیدی در ترکیب با انرژی وزن‌ها است. کانال‌ها یا ابعادی از شبکه که فعال‌سازی آن‌ها در داده‌های کالیبراسیون دارای نوسان بالاتری هستند، توانمندی بیشتری در تمایز و تفکیک توکن‌ها دارا بوده و نباید هرس شوند. این ماژول با پشتیبانی بومی از مدل‌های انکودر و دیکودر، امکان هرس کانال و همچنین تنکی ساختاریافته $2:4$ (\lr{N:M Sparsity}) شتاب‌یافته روی پردازنده‌های گرافیکی مدرن را فراهم می‌سازد.

\subsection{مبانی نظری و فرمول‌بندی ریاضی}
در فاز کالیبراسیون، واریانس مشخصه‌های فعال‌سازی ورودی بر روی دسته‌های کالیبراسیون محاسبه و جمع‌آوری می‌شود. فرض کنید ورودی به لایه خطی $\mathbf{X} \in \mathbb{R}^{B \times d_{\text{in}}}$ باشد. واریانس تجربی هر ویژگی ورودی $j \in \{1, \dots, d_{\text{in}}\}$ به صورت زیر تعریف می‌شود:
\begin{equation}
\mathbf{s}_j = \text{Var}_B(\mathbf{X}_{:, j}) = \frac{1}{B - 1} \sum_{i=1}^B (\mathbf{X}_{i, j} - \bar{\mathbf{X}}_j)^2
\end{equation}
که در آن $\bar{\mathbf{X}}_j$ میانگین فعال‌سازی ویژگی $j$-ام است. در صورتی که چند دسته کالیبراسیون موجود باشد، میانگین متحرک این واریانس‌ها در نظر گرفته می‌شود.

سپس امتیاز اهمیت هر ستون ماتریس وزن $\mathbf{W} \in \mathbb{R}^{d_{\text{out}} \times d_{\text{in}}}$ بر اساس یکی از معیارهای نوسانی زیر سنجیده می‌شود:
\begin{enumerate}
    \item \textbf{معیار واریانس نوسان وزن‌دار (\lr{WIFV - Weighted Input Fluctuation Variance}):}
    \begin{equation}
    \mathbf{I}_j = \mathbf{s}_j \cdot \sum_{i=1}^{d_{\text{out}}} \mathbf{W}_{i, j}^2 = \mathbf{s}_j \cdot \|\mathbf{W}_{:, j}\|_2^2
    \end{equation}
    \item \textbf{معیار نُرم نوسان وزن‌دار (\lr{WIFN - Weighted Input Fluctuation Norm}):}
    \begin{equation}
    \mathbf{I}_j = \sqrt{\mathbf{s}_j} \cdot \sum_{i=1}^{d_{\text{out}}} |\mathbf{W}_{i, j}|
    \end{equation}
    \item \textbf{معیار صرفاً نوسان ورودی (\lr{IFV - Input Fluctuation Variance}):}
    \begin{equation}
    \mathbf{I}_j = \mathbf{s}_j
    \end{equation}
\end{enumerate}

بر اساس میزان تنکی هدف $S \in [0, 1)$، ستون‌هایی که کمترین امتیاز $\mathbf{I}_j$ را دارند ماسک شده یا در حالت تنکی ساختاریافته $2:4$، در هر بلوک ۴تایی از ماتریس، دو عنصری که کمترین ضرب اهمیت را دارند صفر می‌شوند:
\begin{equation}
\mathbf{W}_{\text{pruned}} = \mathbf{W} \odot \mathbf{M}, \quad \mathbf{M} \in \{0, 1\}^{d_{\text{out}} \times d_{\text{in}}}
\end{equation}
خطای نسبی بازسازی فروبنیوس پس از اعمال ماسک ارزیابی می‌گردد:
\begin{equation}
\mathcal{E} = \frac{\|\mathbf{W} - \mathbf{W}_{\text{pruned}}\|_F}{\|\mathbf{W}\|_F}
\end{equation}
در صورتی که $\mathcal{E} > \tau_{\text{max}}$ (آستانه خطای بیشینه)، لایه به وضعیت دست‌نخورده اولیه بازگردانده می‌شود (\lr{Fallback}).

\subsection{معماری و پیاده‌سازی}
کلاس \lr{FLAPCompressor} در فایل \lr{Codes/framework/compression\_engine/flap.py} پیاده‌سازی شده و تحت شناسه‌های \lr{flap} و \lr{flap\_compressor} در \lr{COMPRESSOR\_REGISTRY} ثبت گردیده است. ویژگی‌های معماری و مهندسی این ماژول عبارتند از:
\begin{itemize}
    \item \textbf{سازگاری کامل با \lr{Conv1D} و \lr{Linear}:} پیاده‌سازی با تشخیص خودکار نوع لایه، ترانهاده مناسب وزن‌ها را در لایه‌های \lr{c\_attn} و \lr{c\_proj} در \lr{GPT-2} اعمال کرده و سپس ماسک تنکی را اعمال می‌نماید.
    \item \textbf{پشتیبانی از تنکی سخت‌افزاری $2:4$ (\lr{Structured 2:4 Pruning}):} امکان فعال‌سازی پرچم \lr{structured\_24} برای استفاده از شتاب‌دهنده‌های تانسوری \lr{NVIDIA Ampere} و معماری‌های جدیدتر جهت افزایش توان عملیاتی استنتاج.
    \item \textbf{حفاظت از لایه‌های حساس (\lr{Critical Layer Protection}):} فیلترهای استخراج لایه‌ها از دستکاری طبقه‌بند نهایی (\lr{classifier})، سر مدل زبان (\lr{lm\_head})، نرمال‌سازی لایه‌ها و پروجکشن‌های لورای موجود ممانعت می‌کنند.
    \item \textbf{سازوکار بازگشت ایمن (\lr{Safe Fallback}):} در صورت تخطی خطای بازسازی فروبنیوس از آستانه مجاز (\lr{max\_error=0.35})، لایه بدون تغییر باقی مانده و مشخصات آن در گزارش فشرده‌سازی ثبت می‌شود.
\end{itemize}

\subsection{نتایج تجربی}
ارزیابی ماژول \lr{flap} بر روی مدل‌های انکودر و دیکودر استاندارد بر مبنای داده‌های پایگاه ارزیابی جامع \lr{thesis\_final/master\_verified\_results.json} انجام گرفته است. جدول \ref{tab:benchmark_flap} خلاصه عملکرد ثبت‌شده را نمایش می‌دهد:

\begin{table}[H]
\centering
\small
\caption{نتایج تجربی و اعتبارسنجی بنچ‌مارک فشرده‌ساز \lr{flap} روی مدل‌های پایه}
\label{tab:benchmark_flap}
\begin{tabular}{lccccc}
\toprule
\textbf{مدل پایه} \lr{(Base Model)} & \textbf{دقت} \lr{(Accuracy)} & \textbf{شاخص} \lr{(F1-Score)} & \textbf{حافظه} \lr{(Peak VRAM)} & \textbf{توان} \lr{(Throughput)} & \textbf{وضعیت} \\
\midrule
\lr{DistilBERT} & \lr{90.48\%} & \lr{0.9077} & \lr{311.19 MB} & \lr{164,428.56 tok/s} & \lr{PASS} \\
\lr{RoBERTa} & \lr{88.53\%} & \lr{0.8936} & \lr{321.14 MB} & \lr{52,450.16 tok/s} & \lr{PASS} \\
\lr{GPT-2} & \lr{82.68\%} & \lr{0.8405} & \lr{400.26 MB} & \lr{53,356.34 tok/s} & \lr{PASS} \\
\lr{Qwen} & \lr{60.78\%} & \lr{---} & \lr{---} & \lr{---} & \lr{PASS} \\
\bottomrule
\end{tabular}
\end{table}

نتایج به‌دست‌آمده نشان می‌دهد که رویکرد نوسان فعال‌سازی موفق شده است دقت مدل‌های انکودر نظیر \lr{DistilBERT} را تا سطح کم‌نظیر \lr{90.48\%} حفظ کند، در حالی که توان عملیاتی بسیار بالای \lr{164,428.56 tok/s} و مصرف حافظه ناچیز \lr{311.19 MB} را تثبیت کرده است. در مدل‌های دیکودر مانند \lr{GPT-2} نیز دقت \lr{82.68\%} همراه با توان بالای \lr{53,356.34 tok/s} حاصل گردیده است.

\subsection{محدودیت‌ها و تحلیل شکست}
علیرغم عملکرد برجسته در بنچ‌مارک‌های استاندارد، محدودیت‌های زیر برای \lr{FLAP} قابل شناسایی است:
\begin{itemize}
    \item \textbf{افت دقت در مدل‌های بزرگ علّی (\lr{Qwen}):} در مدل‌های پارامتربالا مانند \lr{Qwen}، افت دقت به سطح \lr{60.78\%} مشاهده می‌شود که ناشی از حساسیت فاحش دیکودرهای خودبازگشت به تنکی ستونی بدون تنظیم مجدد ماتریس‌های کلید-مقدار در لایه‌های عمیق است.
    \item \textbf{وابستگی به توزیع داده‌های کالیبراسیون:} محاسبه دقیق نوسانات نیازمند نمونه‌های کالیبراسیون با تنوع زبانی بالا است؛ در صورتی که داده‌های کالیبراسیون محدود یا دارای بایاس باشند، واریانس ویژگی‌ها گمراه‌کننده شده و منجر به هرس نامناسب کانال‌های مهم اطلاعاتی می‌گردد.
    \item \textbf{نیاز به کامپایلر سخت‌افزاری جهت بهره‌برداری کامل از تنکی:} هرس ستونی ماسک‌شده، در صورت عدم حذف فیزیکی ابعاد ماتریس یا عدم استفاده از هسته‌های بهینه‌شده \lr{Sparse Tensor Core}، لزوماً به کاهش خطی زمان استنتاج خام در تمامی سخت‌افزارها منتج نخواهد شد.
\end{itemize}
